\section{Proofs for the Reference Mechanism and CPT Gradient}\label{app:reference-gradient}

\subsection{Same terminal wealth and path-dependent reference}

\begin{proof}
It is enough to compute the two terminal references and then use strict monotonicity of the deterministic CPT value.

For path I, the first step is a gain, so its interim reference is
\[
 r_{\rm I}^{(1)}=(1-\eta_+)r_0+\eta_+x_{\rm hi}.
\]
The first displayed condition in the proposition makes the second step a loss relative to this interim reference.  Hence
\[
 r_{\rm I}=(1-\eta_-)\{(1-\eta_+)r_0+\eta_+x_{\rm hi}\}+\eta_-x_{\rm f}.
\]
For path II, the first step is a loss and
\[
 r_{\rm II}^{(1)}=(1-\eta_-)r_0+\eta_-x_{\rm lo}.
\]
The second displayed condition makes the move to $x_{\rm f}$ a gain, so
\[
 r_{\rm II}=(1-\eta_+)\{(1-\eta_-)r_0+\eta_-x_{\rm lo}\}+\eta_+x_{\rm f}.
\]
Subtracting cancels the coefficient of $r_0$ and gives \eqref{eq:path-reference-gap}.

It remains to translate a nonzero reference gap into a CPT-value gap.  A deterministic relative outcome $y_0$ has survival function equal to one below its corresponding value magnitude and zero above it.  Since $w_\pm^\eps(0)=0$ and $w_\pm^\eps(1)=1$, its CPT value is
$v_+(y)-v_-(y)$.  The shifted-power gain term is strictly increasing on $y>0$, and $-v_-(y)$ is strictly increasing on $y<0$; the two pieces meet continuously at zero.  Thus the signed deterministic CPT value is strictly increasing in $y$.  If $r_{\rm I}\ne r_{\rm II}$, then $x_{\rm f}-r_{\rm I}\ne x_{\rm f}-r_{\rm II}$, so the two CPT values differ.
\end{proof}

\subsection{Reference contraction and CPT sensitivity}

\begin{proof}
To prove \eqref{eq:reference-value-sensitivity}, it is enough to establish three bounds: contraction of the reference recursion, Lipschitz continuity of the value transformation, and an $L^1$ survival-function coupling inequality.

Fix a wealth value $x$.  The map $r\mapsto\Phi(r,x)$ is continuous, with slope $1-\eta_+$ when $r\le x$ and slope $1-\eta_-$ when $r>x$.  Splitting an interval at $x$ when necessary therefore gives
\[
 \abs{\Phi(r,x)-\Phi(r',x)}
 \le(1-\eta_{\min})\abs{r-r'}.
\]
Under Assumption \ref{ass:reference-coupling}, the two compared trajectories have the same wealth path.  Iterating the preceding inequality for $h$ steps proves \eqref{eq:reference-contraction}.  Since their terminal wealth is common,
\begin{equation}\label{eq:relative-reference-coupling}
 \abs{Y(r)-Y(r')}
 \le(1-\eta_{\min})^h\abs{r-r'}.
\end{equation}

For $x\ge0$,
\[
 \frac{d}{dx}u_{\alpha,\eps_v}(x)
 =\alpha x(x^2+\eps_v^2)^{\alpha/2-1}
 \le\alpha\eps_v^{\alpha-1},
\]
where the bound uses $x\le(x^2+\eps_v^2)^{1/2}$ and $\alpha\le1$.  Hence
\[
 \abs{v_+(y)-v_+(y')}
 +\abs{v_-(y)-v_-(y')}
 \le L_v\abs{y-y'}.
\]
Combining this with \eqref{eq:relative-reference-coupling} controls the sum of the two terminal value-magnitude differences.

Finally, if $V,V'\in[0,B_v]$ are coupled nonnegative random variables, Tonelli's theorem gives
\begin{align*}
 \int_0^{B_v}\abs{P(V>y)-P(V'>y)}\,dy
 &\le \E\int_0^{B_v}
 \abs{\1\{V>y\}-\1\{V'>y\}}\,dy\\
 &=\E\abs{V-V'}.
\end{align*}
Since $\abs{w_\sigma^\eps(p)-w_\sigma^\eps(q)}\le C_1\abs{p-q}$, applying this inequality to the gain and loss integrals and using the preceding value bound proves \eqref{eq:reference-value-sensitivity}.  The decomposition \eqref{eq:variation-decomposition} follows by adding and subtracting $J(P_{k+1},z_{k+1},r_k;\theta)$ and applying the triangle inequality uniformly in $\theta$.
\end{proof}

\subsection{Score representation of the CPT gradient}

\begin{proof}
We prove the gain part; the loss part follows by the same calculation and is subtracted at the end.  To differentiate
\[
 J^+(\theta)=\int_0^{B_v}w_+^\eps(S_\theta^+(y))\,dy,
\]
it is enough to identify $\nabla S_\theta^+(y)$ and justify moving the derivative through the $y$-integral.

For fixed $y$, the likelihood-ratio identity in \Cref{lem:dirichlet-moments} applied to the bounded indicator gives
\begin{equation}\label{eq:survival-score-derivative}
 \nabla S_\theta^+(y)
 =\E\{\1[v_+(Y)>y]G_\theta\}.
\end{equation}
By Cauchy--Schwarz and Assumption \ref{ass:trajectory}, the norm of the right-hand side is at most $\sqrt{C_{G,2}}$ uniformly in $y$ and $\theta$.  Since $(w_+^\eps)'$ is bounded by $C_1$ and the integration interval is finite, this constant is an integrable dominating function.  The dominated convergence theorem therefore gives
\begin{align*}
 \nabla J^+(\theta)
 &=\int_0^{B_v}(w_+^\eps)'(S_\theta^+(y))
 \E\{\1[v_+(Y)>y]G_\theta\}\,dy\\
 &=\E\left[G_\theta
 \int_0^{B_v}(w_+^\eps)'(S_\theta^+(y))\1[v_+(Y)>y]dy\right],
\end{align*}
where Fubini's theorem is applicable because the integrand is dominated by $C_1B_v\norm{G_\theta}$ and $\E\norm{G_\theta}<\infty$.

The score has mean zero.  Therefore subtracting the deterministic quantity $S_\theta^+(y)$ inside the integral does not change the expectation.  This yields the gain component of \eqref{eq:cpt-multiplier}.  Repeating the argument for the loss side and subtracting proves \eqref{eq:score-gradient-identity}.  The frozen-simulator identity follows under the same conditional argument with its trajectory law in place of the true law.
\end{proof}
